\chapter{Analyse et spécification des besoins}
\label{ch:analyse}

\section{Identification des acteurs}

LexHub est un SaaS multi-cabinets. L'unité d'isolation n'est pas l'utilisateur isolé, mais le \textbf{cabinet} : tous les dossiers, pièces, plaidoiries et journaux d'un locataire sont inaccessibles aux autres. Quatre rôles humains s'y articulent. Le tableau~\ref{tab:acteurs} décrit, pour chacun, la responsabilité métier --- ce que la personne vient \emph{faire} dans le produit, et ce qui lui est volontairement refusé.

\begin{table}[H]
\centering
\caption{Acteurs de LexHub et responsabilités métier}
\label{tab:acteurs}
\small
\renewcommand{\arraystretch}{1.28}
\begin{tabular}{@{}>{\raggedright\arraybackslash}p{3.1cm}>{\ttfamily}l>{\raggedright\arraybackslash}p{7.4cm}@{}}
\toprule
\textbf{Acteur} & \textbf{Rôle} & \textbf{Responsabilité dans le produit} \\
\midrule
Cabinet (locataire) & --- &
Espace isolé créé à l'inscription. Porte les données, la charte graphique, l'abonnement et les quotas. N'est pas un utilisateur : c'est le périmètre de confidentialité. \\
\midrule
Avocat & avocat &
Mène l'atelier de plaidoirie (dossier, dictée, correction, rédaction, export), gère ses dossiers et le planning, valide un brouillon pour le rendre visible à l'équipe. Ne paie pas et ne gère pas l'abonnement. \\
\midrule
Assistant juridique & assistant\_juridique &
Consulte uniquement les plaidoiries \textbf{déjà validées} ou \textbf{finalisées}, y dépose des pièces, exporte Word (.docx) et PDF. N'ouvre pas l'atelier de dictée, n'appelle pas les fonctions d'IA de rédaction, de synthèse documentaire ni de fiche orale. \\
\midrule
Admin cabinet & admin\_cabinet &
Dispose de tout le périmètre avocat, plus l'équipe (invitations e-mail), l'abonnement, le logo, le parrainage inter-cabinets et la clôture du cabinet. \\
\midrule
Admin plateforme & admin\_webvoice &
Supervise les cabinets (suspension, purge, plans globaux) sans ouvrir le contenu des affaires --- secret professionnel. \\
\bottomrule
\end{tabular}
\end{table}

\section{Cas d'utilisation}
\label{sec:uc}

La figure~\ref{fig:uc} retient les cas d'utilisation \textbf{spécifiques} à LexHub (hors login / OTP génériques). L'admin cabinet hérite de l'atelier avocat. La source Mermaid (aperçu GitHub / mermaid.live) est fournie dans \texttt{diagrams/uc-lexhub.mmd}.

\begin{figure}[H]
\centering
\begin{tikzpicture}[font=\footnotesize, x=1cm, y=1cm]
  \draw[royal!65, thick, rounded corners=5pt] (2.2,-3.6) rectangle (11.6,4.6);
  \node[font=\bfseries\small,royal] at (6.9,4.25) {LexHub};

  \umlactor{0.5}{2.8}{Avocat};
  \umlactor{0.5}{0.2}{Assistant};
  \umlactor{0.5}{-2.2}{Admin cabinet};
  \umlactor{13.3}{1.0}{Admin plateforme};

  \draw[-{Triangle[open,length=5pt,width=5pt]},thick,royal!70]
    (0.5,-1.55) -- (0.5,2.05) node[midway,left,font=\tiny]{hérite};

  \node[ellipse,draw=royal,thick,align=center,inner sep=2.5pt,fill=white,minimum width=2.55cm] (uc1) at (4.2,3.3) {Dossier\\\& pièces};
  \node[ellipse,draw=royal,thick,align=center,inner sep=2.5pt,fill=white,minimum width=2.55cm] (uc2) at (6.9,3.3) {Dicter /\\audio};
  \node[ellipse,draw=royal,thick,align=center,inner sep=2.5pt,fill=white,minimum width=2.55cm] (uc3) at (9.6,3.3) {Corriger\\forme};
  \node[ellipse,draw=royal,thick,align=center,inner sep=2.5pt,fill=white,minimum width=2.55cm] (uc4) at (4.2,1.6) {Rédiger\\plaidoirie};
  \node[ellipse,draw=royal,thick,align=center,inner sep=2.5pt,fill=white,minimum width=2.55cm] (uc5) at (6.9,1.6) {Valider};
  \node[ellipse,draw=royal,thick,align=center,inner sep=2.5pt,fill=white,minimum width=2.55cm] (uc6) at (9.6,1.6) {Assigner};
  \node[ellipse,draw=royal,thick,align=center,inner sep=2.5pt,fill=white,minimum width=2.55cm] (uc7) at (4.2,-0.1) {Exporter\\Word/PDF};
  \node[ellipse,draw=royal,thick,align=center,inner sep=2.5pt,fill=white,minimum width=2.55cm] (uc8) at (6.9,-0.1) {Voir\\validées};
  \node[ellipse,draw=royal,thick,align=center,inner sep=2.5pt,fill=white,minimum width=2.55cm] (uc9) at (9.6,-0.1) {Inviter\\équipe};
  \node[ellipse,draw=royal,thick,align=center,inner sep=2pt,fill=white,minimum width=2.35cm] (uc10) at (4.2,-1.8) {Abonnement};
  \node[ellipse,draw=royal,thick,align=center,inner sep=2pt,fill=white,minimum width=2.35cm] (uc11) at (6.9,-1.8) {Planning};
  \node[ellipse,draw=royal,thick,align=center,inner sep=2pt,fill=white,minimum width=2.35cm] (uc12) at (9.6,-1.55) {Suspendre\\cabinet};
  \node[ellipse,draw=royal,thick,align=center,inner sep=2pt,fill=white,minimum width=2.35cm] (uc12a) at (9.6,-3.0) {Usage\\\& plans};

  \draw[dashed,-{Stealth},thick,royal!60] (uc2) -- node[above,font=\tiny]{<<i>>} (uc3);
  \draw[dashed,-{Stealth},thick,royal!60] (uc6) -- node[above,font=\tiny]{<<i>>} (uc5);

  \draw (1.2,2.6) -- (uc1.west);
  \draw (1.2,2.6) -- (uc2.west);
  \draw (1.2,2.6) -- (uc3.west);
  \draw (1.2,2.6) -- (uc4.west);
  \draw (1.2,2.6) -- (uc5.west);
  \draw (1.2,2.6) -- (uc6.west);
  \draw (1.2,2.6) -- (uc7.west);
  \draw (1.2,2.6) -- (uc11.west);

  \draw (1.2,0.0) -- (uc8.west);
  \draw (1.2,0.0) -- (uc7.south west);

  \draw (1.2,-2.4) -- (uc9.west);
  \draw (1.2,-2.4) -- (uc10.west);

  \draw (12.6,0.8) -- (uc12.east);
  \draw (12.6,0.8) -- (uc12a.east);
\end{tikzpicture}
\caption{Diagramme de cas d'utilisation UML --- LexHub (cas principaux)}
\label{fig:uc}
\end{figure}

\section{Contrôle d'accès (RBAC) --- vérité produit}
\label{sec:acces}

Le « diagramme d'accès » de LexHub n'est pas un schéma réseau : c'est la \textbf{matrice des droits} réellement appliquée par l'API (JWT + rôle + \texttt{cabinet\_id}). La figure~\ref{fig:acces} synthétise ce qui est autorisé ; tout le reste est refusé côté serveur (pas seulement masqué dans l'UI).

\begin{table}[H]
\centering
\caption{Matrice de contrôle d'accès LexHub (production)}
\label{fig:acces}
\small
\renewcommand{\arraystretch}{1.22}
\begin{tabular}{@{}p{4.4cm}cccc@{}}
\toprule
\textbf{Capacité} & \textbf{Avocat} & \textbf{Assistant} & \textbf{Admin cab.} & \textbf{Admin plat.} \\
\midrule
Isolation par cabinet & oui & oui & oui & n/a \\
Atelier dictée / rédaction IA & oui & non & oui & non \\
Voir brouillons (\texttt{draft}) & oui & non & oui & non \\
Voir plaidoiries validées / finalisées & oui & oui & oui & non \\
Exporter Word (.docx) / PDF & oui & oui\textsuperscript{*} & oui & non \\
Valider / assigner & oui & non & oui & non \\
Inviter l'équipe & non & non & oui & non \\
Abonnement / packs / logo & non & non & oui & non \\
Suspendre / purger un cabinet & non & non & non & oui \\
Plans globaux, leads, logs, coach & non & non & non & oui \\
Lire le contenu des affaires & oui\textsuperscript{\dag} & limité\textsuperscript{*} & oui\textsuperscript{\dag} & \textbf{non} \\
\bottomrule
\end{tabular}

\vspace{0.35em}
{\footnotesize
\textsuperscript{*}~uniquement les plaidoiries au statut \textit{validated} ou \textit{finalized} (pas les brouillons).
\textsuperscript{\dag}~uniquement dans \emph{son} cabinet (\texttt{cabinet\_id}).
}
\end{table}

\paragraph{Règles à retenir pour la soutenance.}
\begin{itemize}[leftmargin=1.4em]
  \item L'\textbf{assignation} oriente le travail ; elle \textbf{ne crée pas} un droit exclusif. Tout assistant du cabinet voit les plaidoiries validées.
  \item L'\textbf{admin plateforme} gère l'exploitation (cabinets, plans, audit) \textbf{sans} ouvrir dossiers ni pièces --- secret professionnel.
  \item Après trois échecs de login : Cloudflare Turnstile ; le jeton est vérifié par l'API.
\end{itemize}

\section{Exigences fonctionnelles}

Les exigences ci-dessous décrivent le \textbf{produit livré}. Elles ne portent pas de priorité de backlog : le rapport documente un système déjà en production, pas un arbitrage de sprint.

\paragraph{Espace cabinet et accès.}
\begin{itemize}[leftmargin=1.4em]
  \item \textbf{EF-CAB-1.} Inscription d'un cabinet : vérification de l'e-mail par OTP, puis accès immédiat (cabinet créé actif, plan Découverte). L'admin plateforme peut ensuite suspendre ou réactiver un cabinet ; il n'y a pas de file d'approbation à l'inscription.
  \item \textbf{EF-CAB-2.} Connexion e-mail / mot de passe, session authentifiée, interface bilingue français / arabe. Après trois échecs de login, un défi anti-robot Cloudflare Turnstile est exigé ; le jeton est vérifié côté API. L'inscription n'affiche pas ce widget (preuve d'e-mail par OTP).
  \item \textbf{EF-CAB-3.} Isolation stricte des données par cabinet.
\end{itemize}

\paragraph{Atelier de plaidoirie.}
\begin{itemize}[leftmargin=1.4em]
  \item \textbf{EF-PL-1.} Parcours en cinq étapes : \textbf{Dossier \& pièces}, \textbf{Dictée}, \textbf{Transcription}, \textbf{Rédaction IA}, \textbf{Export}. La feuille de clarification et le chat de révision sont \emph{intégrés} aux étapes 3--4 ; ce ne sont pas des étapes supplémentaires du rail.
  \item \textbf{EF-PL-2.} Dictée : reconnaissance vocale native du navigateur (Chrome / Edge, Web Speech API) avec aperçu instantané, ou capture micro via \textbf{MediaRecorder $\rightarrow$ Whisper} (OpenAI) --- notamment sous Firefox --- ou import fichier audio. La langue de dictée (arabe / français) peut être détectée \textbf{automatiquement} sans basculer toute l'interface.
  \item \textbf{EF-PL-3.} Bouton \textbf{Corriger} : correction de forme (arabe standard), tolérante aux fautes de frappe / STT, sans résumé ni invention ; surlignage des substitutions à relire ; alertes lorsque des faits contradictoires (numéros, dates, montants, noms) coexistent dans la dictée.
  \item \textbf{EF-PL-4.} Composition de la plaidoirie à partir du dossier, des pièces et de la dictée corrigée, sans inventer un fait, un article ou un montant absent de ces sources. Avant rédaction : \textbf{clarification QCM} (0 à 5 questions) générée par le LLM après lecture anonymisée du dossier et de la dictée ; l'avocat confirme ou passe.
  \item \textbf{EF-PL-5.} Après rédaction : panneau \textbf{Assistant IA --- modifications} (chat) : question et/ou consigne en langage naturel (fautes acceptées) ; l'IA \textbf{répond} dans le chat \textbf{et} met à jour le document (ex.\ remplir \texttt{خامساً} / \texttt{لهذه الأسباب}). Contrôles : \textbf{Annuler} / \textbf{Modifier la requête}, micro avec bandeau d'écoute et bouton \textbf{Terminé} (détection FR/AR).
  \item \textbf{EF-PL-6.} Export officiel Word (.docx) --- standard modifiable des tribunaux marocains (en-tête bilingue, Noto Naskh RTL, armoiries) --- \textbf{et} export PDF (même mise en page). Après export réussi : statut \textbf{finalisée} (mise à jour côté client).
\end{itemize}

\paragraph{Collaboration.}
\begin{itemize}[leftmargin=1.4em]
  \item \textbf{EF-COL-1.} Un brouillon reste invisible à l'assistant tant que son statut n'est pas \textbf{validé} (ou \textbf{finalisé}). Tous les assistants du cabinet voient alors cette plaidoirie.
  \item \textbf{EF-COL-2.} L'action \textbf{Assigner} rattache un assistant, valide le brouillon s'il ne l'était pas, et met la plaidoirie en avant sur son tableau de bord. Ce n'est pas le filtre d'autorisation : le filtre est le statut validé / finalisé.
  \item \textbf{EF-COL-3.} L'admin cabinet invite avocats et assistants par \textbf{e-mail}. L'invité ouvre le lien, choisit son mot de passe, et accède au cabinet. Il n'existe plus de code cabinet partagé.
  \item \textbf{EF-COL-4.} Trois statuts : \texttt{draft} (brouillon), \texttt{validated} (partage assistant), \texttt{finalized} (exportée).
\end{itemize}

\paragraph{Dossier, documents, planning.}
\begin{itemize}[leftmargin=1.4em]
  \item \textbf{EF-DOS-1.} Dossiers (domaine, identité judiciaire). Pièces : OCR cloud, synthèse et questions, réservées à l'avocat et à l'admin cabinet.
  \item \textbf{EF-DOS-2.} Fiche orale d'audience (\textbf{clé de barre}) : document de préparation généré depuis le dossier, distinct de la plaidoirie écrite.
  \item \textbf{EF-DOS-3.} Planning d'audiences et de délais, avec rappels.
\end{itemize}

\paragraph{Pilotage et économie.}
\begin{itemize}[leftmargin=1.4em]
  \item \textbf{EF-ECO-1.} Abonnement, essai hybride, packs de crédits, sièges supplémentaires --- chapitre~\ref{ch:economie}.
  \item \textbf{EF-ECO-2.} Parrainage d'un cabinet par un autre, via un lien d'inscription (le code n'ouvre aucun dossier).
  \item \textbf{EF-ECO-3.} Journal d'activité, notifications, quitter / supprimer son compte (sauf admin plateforme).
\end{itemize}

\section{Exigences non fonctionnelles}

\begin{table}[H]
\centering
\caption{Exigences non fonctionnelles}
\label{tab:enf}
\small
\renewcommand{\arraystretch}{1.2}
\begin{tabular}{@{}clp{9.8cm}@{}}
\toprule
\textbf{ID} & \textbf{Catégorie} & \textbf{Exigence} \\
\midrule
ENF-1 & Confidentialité & Isolation par \texttt{cabinet\_id} ; HTTPS / TLS ; mots de passe robustes ; sessions JWT. Le navigateur n'ouvre aucune session SQL. \\
ENF-2 & Audit & Journal des actions sensibles (authentification, validation, assignation, administration). \\
ENF-3 & Disponibilité & Production sur VPS, reverse-proxy Caddy, conteneurs Docker. \\
ENF-4 & Usabilité & Interface responsive ; parcours plaidoirie guidé ; bilingue FR / AR. \\
ENF-5 & Intégrité IA & Pas d'invention de fond ; contrôle déterministe des contradictions factuelles. \\
ENF-6 & Résistance & Limitation de débit sur l'authentification et les endpoints sensibles ; Cloudflare Turnstile après trois échecs de connexion (production). \\
ENF-7 & Observabilité & E-mails transactionnels ; notifications temps réel ; page Activité. \\
\bottomrule
\end{tabular}
\end{table}

\section{Règles de gestion}

\begin{enumerate}[leftmargin=1.4em]
  \item Un brouillon n'est pas visible des assistants.
  \item La \textbf{validation} rend la plaidoirie visible de \emph{tous} les assistants du cabinet.
  \item L'\textbf{export} Word (.docx) / PDF, puis statut \textbf{finalisée} via l'UI (distinct de la validation assistant).
  \item L'\textbf{assignation} valide si besoin et met la plaidoirie en avant pour un assistant nommé ; elle ne crée pas un droit d'accès exclusif.
  \item L'\textbf{épinglage} ordonne la liste, sans changer le statut.
  \item L'admin plateforme suspend un cabinet ; il ne lit pas les dossiers.
\end{enumerate}

\section{Périmètre fonctionnel}
\label{sec:perimetre}

\paragraph{Dans le périmètre (livré).}
Atelier en cinq étapes ; isolation multi-cabinets et quatre rôles ; invitation d'équipe ; export Word/PDF ; abonnement (essai, plans, packs, sièges) ; parrainage ; supervision plateforme sans lecture des affaires.

\paragraph{Hors périmètre (volontaire).}
Pas de portail ministériel, ni de comptabilité cabinet, ni de chatbot juridique autonome. Le paiement carte opérateur dispose d'un environnement technique prêt ; le débit réel n'est pas activé à la date du rapport (simulation).
